% !TeX root = ../main.tex
\section{Limitations}
\label{sec:limitations}

The main limitation lies in the source material. The rules of Cờ Trận were
recorded from the recollection of \emph{one} person, the author, many years
after they stopped playing (Section~\ref{sec:provenance}). We have not found
a second person to cross-check them, nor any written document. The record
may therefore differ from how the game was played at other times, by other
groups of children in the village, or in neighbouring villages, and errors of
memory cannot be ruled out. We have no information about the longer history
of the game: when it arose, who devised it, or whether it is related to any
other game.

The software does not yet have an online mode for remote multiplayer play.
The AI's evaluation function uses only two features (piece count and
mobility); deeper strategic factors such as centre control, multi-direction
threats, or positional value are not modelled.

The formal mathematical model in Appendix~\ref{app:formal} is complete only
for the 2-player mode. The extension to 3 and 4 players is specified only at
the configuration level (board size, starting positions---%
Appendix~\ref{app:multistart}) and has not been analysed theoretically,
especially the balance properties of multiplayer games. Notably, the
starting configuration of the 3-player mode is \emph{unbalanced}: two sides
have 13 pieces while the third has 14. We record this configuration exactly
as in the folk game and do \emph{not} alter it, since changing the rules
would falsify the record. Any more balanced version would have to be
presented clearly as a new variant; automated balance measurement and design
methods~\citep{hommarks2007balanced} offer one way to explore this.

The bot-vs-bot tournament in Section~\ref{sec:eval} has so far only compared
search depths in the 2-player mode. The first-mover advantage measurement
(Section~\ref{sec:eval_tournament}) is also only in the 2-player mode and at
one fixed depth ($d{=}2$, where the second mover appears to have an
advantage); there is no quantitative study yet of the balance among colours
in the 3- and 4-player modes---the questions of whether moving first or
second confers a systematic advantage at other depths, and whether an
implicit coalition strategy is stable, remain open.

As for the terminal condition, Cờ Trận has \emph{no} move limit and no
repetition-based draw rule, in any mode. In theory a game could go on forever
without any side being blockaded. This is a property of the folk game itself:
in real play, a four-player game could last three to four hours
(Section~\ref{sec:provenance}). The software keeps it that way and sets no
move limit. The AI's anti-repetition mechanism (Section~\ref{ssec:pitfalls})
makes endless repetition rare when the machine plays, but the rules do not
rule it out.
